% GENERATED FILE. Edit canonical lessons, then run npm run generate:books.
% canonical-source-hash: fnv1a64:38dfa287558f6eba
% canonical-lessons: JA-C212-kion, JA-C212-agaru, JA-C212-hieru, JA-C212-kooru, JA-C212-tokeru

\chapter{Cold Nights and Warm Afternoons}
\label{ch:ja-cold-nights-and-warm-afternoons}
\hlchaptermodality{212}

\begin{chapteropening}
\textbf{By the end of this chapter:} \emph{I can say how hot or cold it will get.}

Complete the last lesson of chapter 212: I can say how hot or cold it will get.
\end{chapteropening}

\section[\texorpdfstring{\ja{きおん} (kion)}{きおん (kion)}]{\ja{きおん} (kion) — air temperature}

\label{lesson:JA-C212-kion}



\begin{quote}
Before the new one: say the Japanese for a small bird, then the Japanese for to increase.
\end{quote}

\subsection*{You'll want to know: \ja{きおん}}
\textbf{\ja{きおん}} — \emph{kion} — \textquotedblleft{}air temperature\textquotedblright{}.

\begin{quote}
\textbf{\ja{あしたは} \ja{きおんが} \ja{たかいでしょう。}} — \emph{ashita wa kion ga takai deshō} — \textquotedblleft{}The temperature will probably be high tomorrow.\textquotedblright{}
\end{quote}

\begin{cousinweb}
\textbf{\ja{きおん}} (\emph{kion}) is Sino-Japanese (kango): two Chinese characters read with their on'yomi, readings borrowed from Chinese. They are \emph{ki}, \textquotedblleft{}air\textquotedblright{} (Middle Chinese \emph{khj+jH}) and \emph{on}, \textquotedblleft{}warm\textquotedblright{} (\emph{'won}) (KANJIDIC2; Middle Chinese in Baxter's transcription).
\end{cousinweb}

\begin{grammarlens}[title={What you've built}]
One more word for everyday life at A2.
\end{grammarlens}

\subsection*{Your turn}
\begin{itemize}
\raggedright
  \item \emph{Say it:} \emph{kion}
  \item \emph{Say it:} \emph{kion}, once more
  \item \emph{Say it:} \emph{ashita wa kion ga takai deshō} — The temperature will probably be high tomorrow.
  \item \emph{Recall:} say the Japanese for a small bird, then the Japanese for to increase, then say \emph{kion} again
\end{itemize}

\begin{tcolorbox}[breakable,colback=teal!4,colframe=teal!35!black,title={Before you move on}]
What does \ja{きおん} mean? (\textquotedblleft{}Air temperature\textquotedblright{}.) Say it once more.
\end{tcolorbox}
\section[\texorpdfstring{\ja{あがる} (agaru)}{あがる (agaru)}]{\ja{あがる} (agaru) — to rise; to go up}

\label{lesson:JA-C212-agaru}



\begin{quote}
Before the new one: say the Japanese for to increase, then the Japanese for air temperature.
\end{quote}

\subsection*{You'll want to know: \ja{あがる}}
\textbf{\ja{あがる}} — \emph{agaru} — \textquotedblleft{}to rise; to go up\textquotedblright{}.

\begin{quote}
\textbf{\ja{あした} \ja{きおんが} \ja{あがります。}} — \emph{ashita kion ga agarimasu} — \textquotedblleft{}The temperature will rise tomorrow.\textquotedblright{}
\end{quote}

Polite \textbf{\ja{あがります}}, \emph{agarimasu}; a godan verb.

\begin{cousinweb}
\textbf{\ja{あがる}} (\emph{agaru}) is a native Japanese word (wago), inherited rather than borrowed from Chinese; its kanji was attached later as a kun'yomi, and its older origin is not settled (JMdict, KANJIDIC2).
\end{cousinweb}

\begin{grammarlens}[title={What you've built}]
One more word for everyday life at A2.
\end{grammarlens}

\subsection*{Your turn}
\begin{itemize}
\raggedright
  \item \emph{Say it:} \emph{agaru}
  \item \emph{Say it:} \emph{agaru}, once more
  \item \emph{Say it:} \emph{ashita kion ga agarimasu} — The temperature will rise tomorrow.
  \item \emph{Recall:} say the Japanese for to increase, then the Japanese for air temperature, then say \emph{agaru} again
\end{itemize}

\begin{tcolorbox}[breakable,colback=teal!4,colframe=teal!35!black,title={Before you move on}]
What does \ja{あがる} mean? (\textquotedblleft{}To rise; to go up\textquotedblright{}.) Say it once more.
\end{tcolorbox}
\section[\texorpdfstring{\ja{ひえる} (hieru)}{ひえる (hieru)}]{\ja{ひえる} (hieru) — to get cold; to grow chilly}

\label{lesson:JA-C212-hieru}



\begin{quote}
Before the new one: say the Japanese for air temperature, then the Japanese for to rise.
\end{quote}

\subsection*{You'll want to know: \ja{ひえる}}
\textbf{\ja{ひえる}} — \emph{hieru} — \textquotedblleft{}to get cold; to grow chilly\textquotedblright{}.

\begin{quote}
\textbf{\ja{こんやは} \ja{ひえます。}} — \emph{kon'ya wa hiemasu} — \textquotedblleft{}It will get cold tonight.\textquotedblright{}
\end{quote}

Polite \textbf{\ja{ひえます}}, \emph{hiemasu}; an ichidan verb.

\begin{cousinweb}
\textbf{\ja{ひえる}} (\emph{hieru}) is a native Japanese word (wago), inherited rather than borrowed from Chinese; its kanji was attached later as a kun'yomi, and its older origin is not settled (JMdict, KANJIDIC2).
\end{cousinweb}

\begin{grammarlens}[title={What you've built}]
One more word for everyday life at A2.
\end{grammarlens}

\subsection*{Your turn}
\begin{itemize}
\raggedright
  \item \emph{Say it:} \emph{hieru}
  \item \emph{Say it:} \emph{hieru}, once more
  \item \emph{Say it:} \emph{kon'ya wa hiemasu} — It will get cold tonight.
  \item \emph{Recall:} say the Japanese for air temperature, then the Japanese for to rise, then say \emph{hieru} again
\end{itemize}

\begin{tcolorbox}[breakable,colback=teal!4,colframe=teal!35!black,title={Before you move on}]
What does \ja{ひえる} mean? (\textquotedblleft{}To get cold; to grow chilly\textquotedblright{}.) Say it once more.
\end{tcolorbox}
\section[\texorpdfstring{\ja{こおる} (kōru)}{こおる (kōru)}]{\ja{こおる} (kōru) — to freeze}

\label{lesson:JA-C212-kooru}



\begin{quote}
Before the new one: say the Japanese for to rise, then the Japanese for to get cold.
\end{quote}

\subsection*{You'll want to know: \ja{こおる}}
\textbf{\ja{こおる}} — \emph{kōru} — \textquotedblleft{}to freeze\textquotedblright{}.

\begin{quote}
\textbf{\ja{あしたの} \ja{あさ} \ja{いけが} \ja{こおります。}} — \emph{ashita no asa ike ga kōrimasu} — \textquotedblleft{}The pond will freeze tomorrow morning.\textquotedblright{}
\end{quote}

Polite \textbf{\ja{こおります}}, \emph{kōrimasu}; a godan verb. The noun \textbf{\ja{こおり}}, \emph{kōri}, means \textquotedblleft{}ice\textquotedblright{}.

\begin{cousinweb}
\textbf{\ja{こおる}} (\emph{kōru}) is a native Japanese word (wago), inherited rather than borrowed from Chinese; its kanji was attached later as a kun'yomi, and its older origin is not settled (JMdict, KANJIDIC2).
\end{cousinweb}

\begin{grammarlens}[title={What you've built}]
One more word for everyday life at A2.
\end{grammarlens}

\subsection*{Your turn}
\begin{itemize}
\raggedright
  \item \emph{Say it:} \emph{kōru}
  \item \emph{Say it:} \emph{kōru}, once more
  \item \emph{Say it:} \emph{ashita no asa ike ga kōrimasu} — The pond will freeze tomorrow morning.
  \item \emph{Recall:} say the Japanese for to rise, then the Japanese for to get cold, then say \emph{kōru} again
\end{itemize}

\begin{tcolorbox}[breakable,colback=teal!4,colframe=teal!35!black,title={Before you move on}]
What does \ja{こおる} mean? (\textquotedblleft{}To freeze\textquotedblright{}.) Say it once more.
\end{tcolorbox}
\section[\texorpdfstring{\ja{とける} (tokeru)}{とける (tokeru)}]{\ja{とける} (tokeru) — to melt; to thaw}

\label{lesson:JA-C212-tokeru}



\begin{quote}
Before the new one: say the Japanese for to get cold, then the Japanese for to freeze.
\end{quote}

\subsection*{You'll want to know: \ja{とける}}
\textbf{\ja{とける}} — \emph{tokeru} — \textquotedblleft{}to melt; to thaw\textquotedblright{}.

\begin{quote}
\textbf{\ja{らいしゅう} \ja{ゆきが} \ja{とけます。}} — \emph{raishū yuki ga tokemasu} — \textquotedblleft{}The snow will melt next week.\textquotedblright{}
\end{quote}

Polite \textbf{\ja{とけます}}, \emph{tokemasu}; an ichidan verb.

\begin{cousinweb}
\textbf{\ja{とける}} (\emph{tokeru}) is a native Japanese word (wago), inherited rather than borrowed from Chinese; its kanji was attached later as a kun'yomi, and its older origin is not settled (JMdict, KANJIDIC2).
\end{cousinweb}

\begin{grammarlens}[title={What you've built}]
One more word for everyday life at A2.
\end{grammarlens}

\subsection*{Your turn}
\begin{itemize}
\raggedright
  \item \emph{Say it:} \emph{tokeru}
  \item \emph{Say it:} \emph{tokeru}, once more
  \item \emph{Say it:} \emph{raishū yuki ga tokemasu} — The snow will melt next week.
  \item \emph{Recall:} say the Japanese for to get cold, then the Japanese for to freeze, then say \emph{tokeru} again
\end{itemize}

\begin{tcolorbox}[breakable,colback=teal!4,colframe=teal!35!black,title={Before you move on}]
What does \ja{とける} mean? (\textquotedblleft{}To melt; to thaw\textquotedblright{}.) Say it once more.
\end{tcolorbox}
